\chapter{Modern Video Compression Standards}

\section{Table of Contents}

\begin{itemize}
  \item Core Idea
  \item Key Concepts
  \item Theory and Formulas
  \item Visual Schemes
  \item Examples
\end{itemize}

\section{Core Idea}

Modern video standards from MPEG-1 to H.266/VVC keep the same \textbf{hybrid video coding} paradigm:

\begin{itemize}
  \item predict current block from spatial or temporal neighbors;
  \item code only prediction residual;
  \item transform and quantize residual;
  \item entropy-code coefficients and side information;
  \item reconstruct frame inside encoder to keep decoder and encoder synchronized.
\end{itemize}

\begin{calloutbox}{red!70!black}{Universal hybrid paradigm}
\textbf{Statement:} modern video standards share prediction, residual coding, transform, quantization, entropy coding, in-loop reconstruction, and reference-frame buffering.

\textbf{Meaning:} codec generations mainly differ in how much freedom they give to partitioning, prediction, transforms, filters, and transport syntax.
\end{calloutbox}

\begin{center}\includegraphics[width=\linewidth]{ch09/universal-hybrid-encoder.png}\end{center}
Universal hybrid encoder: mode decision chooses Intra or Inter prediction, residual is transform-coded, and decoded reconstruction re-enters frame buffer to prevent drift.

\section{Key Concepts}

\subsection{What Standards Define}

Video standards do \textbf{not} standardize full encoder optimization. They standardize:

\begin{itemize}
  \item bitstream syntax;
  \item decoder behavior;
  \item conformance rules.
\end{itemize}

Encoder choices such as motion search, rate control, partition pruning, and perceptual tuning remain open.

\begin{center}\includegraphics[width=\linewidth]{ch09/standard-scope-decoder.png}\end{center}
Standard scope: decoder and bitstream are fixed; encoder implementation is competitive design space.

\begin{calloutbox}{red!70!black}{Interoperability}
Any compliant decoder must produce same reconstruction from same valid bitstream. Encoders may differ in RDO search, heuristics, machine-learning pruning, and hardware scheduling.
\end{calloutbox}

\subsection{Historical Evolution}

Generations improve compression by increasing degrees of freedom:

\begin{center}\fitwidth{%
\begin{tabular}{|l|l|l|}
\hline
\textbf{Standard} & \textbf{Main design step} & \textbf{Cost} \\ \hline
\textbf{MPEG-2} & basic hybrid coding & low complexity, weak compression \\ \hline
\textbf{H.264/AVC} & macroblocks, strong ME, CABAC/CAVLC & much better efficiency \\ \hline
\textbf{H.265/HEVC} & CTUs, quad-tree, larger transforms, SAO & higher encoder complexity \\ \hline
\textbf{H.266/VVC} & QT+MTT, affine tools, MTS, ALF & very high complexity \\ \hline
\textbf{AV1} & royalty-free web codec, superblocks, strong filters & high complexity, broad web focus \\ \hline
\end{tabular}}
\end{center}

\begin{center}\includegraphics[width=0.833\linewidth]{ch09/codec-evolution-timeline.png}\end{center}
Codec evolution trades lower bitrate for much larger encoder search space.

\begin{calloutbox}{red!70!black}{Golden rule}
New codec generations often target roughly half bitrate for same perceptual quality, but encoder complexity can increase by about $10x$.
\end{calloutbox}

\subsection{Source 2026 Adoption Scenarios}

Course notes frame codec adoption as ecosystem-driven:

\begin{center}\fitwidth{%
\begin{tabular}{|l|l|l|}
\hline
\textbf{Ecosystem} & \textbf{Typical codecs} & \textbf{Main reason} \\ \hline
real-time / legacy & H.264/AVC & universal hardware, low latency \\ \hline
broadcast / mobile & HEVC, H.264 & mature 4K hardware support \\ \hline
web streaming & AV1, H.264, HEVC & AV1 royalty-free appeal, H.264 compatibility \\ \hline
next-gen 8K / VR & VVC, AV1 & maximum compression, higher complexity \\ \hline
\end{tabular}}
\end{center}

Adoption depends on licensing, decoder hardware, battery cost, latency, and platform support, not only compression efficiency.

\subsection{Core Standard Tools}

\begin{center}\fitwidth{%
\begin{tabular}{|l|l|l|}
\hline
\textbf{Layer} & \textbf{Tool} & \textbf{Purpose} \\ \hline
partitioning & MB, CTU, superblock, QT, MTT & adapt block geometry to content \\ \hline
prediction & Intra, Inter, Merge, affine & remove spatial/temporal redundancy \\ \hline
residual coding & DCT/DST/MTS transforms & compact residual energy \\ \hline
quantization & QP-controlled scalar quantization & rate-distortion control \\ \hline
entropy coding & CABAC / arithmetic coding & lossless syntax compression \\ \hline
filtering & deblocking, SAO, ALF, CDEF & avoid artifacts in reference frames \\ \hline
transport & VCL/NAL, NALUs, OBUs & map compressed video to networks/files \\ \hline
parallelism & slices, tiles, WPP & error containment and hardware speed \\ \hline
\end{tabular}}
\end{center}

\section{Theory and Formulas}

\subsection{Rate-Distortion Optimization}

For each coding decision:

\[
J=D+\lambda R
\]

where $D$ is distortion, $R$ is rate, and $\lambda$ controls quality-rate balance.

\begin{calloutbox}{red!70!black}{RDO}
Encoder should not select lowest distortion or lowest rate alone. It selects option minimizing weighted cost $D+\lambda R$.
\end{calloutbox}

Block splitting is also a mode decision:

\[
J_{\text{split}}=\sum_i J_{\text{subblock}_i}
\]

Split if:

\[
J_{\text{split}}<J_B
\]

This principle drives modern partition trees.

\subsection{Space Partitioning Evolution}

\subsubsection{H.264/AVC: Macroblocks}

H.264 uses $16 \times 16$ \textbf{macroblocks}. Variable block-size partitioning supports:

\begin{itemize}
  \item $16 \times 16$;
  \item $16 \times 8$;
  \item $8 \times 16$;
  \item $8 \times 8$, with sub-partitions down to $4 \times 4$.
\end{itemize}

Limitation: $16 \times 16$ maximum is inefficient for HD/4K smooth regions, where larger areas could be predicted with fewer bits.

\subsubsection{H.265/HEVC: CTU and Quad-Tree}

HEVC replaces macroblocks with \textbf{Coding Tree Units (CTUs)} up to $64 \times 64$. CTU splits recursively into square \textbf{Coding Units (CUs)} through quad-tree partitioning.

\begin{center}\includegraphics[width=0.833\linewidth]{ch09/hevc-ctu-quadtree.png}\end{center}
HEVC quad-tree supports large blocks for smooth areas and small blocks around detail or object boundaries.

HEVC separates:

\begin{itemize}
  \item \textbf{CU}: coding decision unit;
  \item \textbf{PU}: prediction unit;
  \item \textbf{TU}: transform unit.
\end{itemize}

\subsubsection{H.266/VVC: QT + MTT}

VVC expands CTUs up to $128 \times 128$ and introduces \textbf{Multi-Type Tree (MTT)}:

\begin{itemize}
  \item quad-tree split;
  \item binary horizontal/vertical split;
  \item ternary horizontal/vertical split with $1:2:1$ ratio.
\end{itemize}

\begin{center}\includegraphics[width=0.833\linewidth]{ch09/vvc-qt-mtt-partitioning.png}\end{center}
QT+MTT gives rectangular blocks, useful for object edges and directional textures, but increases RDO search complexity.

\subsubsection{AV1: Superblock Partitioning}

AV1 uses up to $128 \times 128$ \textbf{superblocks} and 10-way partitioning, including elongated $4:1$ and $1:4$ patterns.

\begin{center}\includegraphics[width=0.833\linewidth]{ch09/av1-superblock-partitioning.png}\end{center}
AV1 partitioning gives high geometric freedom, especially for web-streaming efficiency.

\subsection{RDO Search Burden}

Modern partitioning forms huge decision trees. Reference encoders often use DFS:

\begin{enumerate}
  \item test parent block;
  \item test splits;
  \item evaluate Intra/Inter modes inside each node;
  \item reconstruct blocks in causal order;
  \item compare child cost against parent cost.
\end{enumerate}

Search is partly sequential because top/left reconstructed pixels are needed for prediction.

Production encoders prune using:

\begin{itemize}
  \item texture variance and gradients;
  \item temporal correlation from co-located blocks;
  \item early termination;
  \item learned split classifiers such as LightGBM or small CNNs.
\end{itemize}

\subsection{Intra Prediction}

\textbf{Intra prediction} uses already decoded top and left neighbors to predict current block.

\begin{center}\includegraphics[width=0.7\linewidth]{ch09/intra-prediction-neighbors.png}\end{center}
Spatial prediction is causal: only already reconstructed neighboring samples can be used.

Mode evolution:

\begin{center}\fitwidth{%
\begin{tabular}{|l|l|}
\hline
\textbf{Standard} & \textbf{Intra modes} \\ \hline
H.264/AVC & 4 modes for $16 \times 16$, 9 modes for $4 \times 4$ \\ \hline
H.265/HEVC & 35 modes: 33 directions plus DC and Planar \\ \hline
H.266/VVC & 65+ directional modes plus wide-angle modes \\ \hline
\end{tabular}}
\end{center}

Mode freedom grows with signaling bits:

\[
P \propto 2^B
\]

\begin{center}\includegraphics[width=0.7\linewidth]{ch09/intra-prediction-directional-example.png}\end{center}
Directional modes improve texture prediction but require more signaling.

\begin{center}\includegraphics[width=0.7\linewidth]{ch09/intra-small-block-combinatorics.png}\end{center}
For many small blocks, Intra-mode combinations explode; practical encoders search sequentially and prune.

\subsection{Most Probable Mode}

If VVC exposes 67 Intra modes, naive signaling needs around 7 bits per block. For tiny $4 \times 4$ blocks this can erase prediction gains.

\textbf{Most Probable Mode (MPM)} solves this:

\begin{itemize}
  \item build candidate list from top and left Intra modes;
  \item if selected mode is likely, send short index;
  \item HEVC uses 3 MPM candidates;
  \item VVC uses 6 MPM candidates.
\end{itemize}

\subsection{Inter Prediction}

Inter prediction uses ME/MC and reference frames in the \textbf{Decoded Picture Buffer (DPB)}.

Modern tools:

\begin{itemize}
  \item multiple reference frames;
  \item one P-list or two B-lists;
  \item generalized B-slices;
  \item weighted averaging of predictors;
  \item fractional-pel motion compensation.
\end{itemize}

\begin{center}\includegraphics[width=0.867\linewidth]{ch09/multiple-reference-frames.png}\end{center}
Multiple reference frames let encoder choose best temporal predictor, not only nearest frame.

Sub-pixel precision:

\begin{center}\fitwidth{%
\begin{tabular}{|l|l|}
\hline
\textbf{Codec family} & \textbf{Motion precision} \\ \hline
H.264/HEVC/VVC & quarter-pel, $1/4$ pixel \\ \hline
AV1 & eighth-pel, $1/8$ pixel \\ \hline
\end{tabular}}
\end{center}

\subsection{Motion-Vector Prediction and Merge}

Motion vectors must also be coded efficiently.

\begin{center}\fitwidth{%
\begin{tabular}{|l|l|}
\hline
\textbf{Tool} & \textbf{Idea} \\ \hline
H.264 Skip/Direct & use median of spatial neighbors; if correct, send no MV \\ \hline
HEVC/VVC Merge & build candidate list from spatial and temporal candidates; send index \\ \hline
\end{tabular}}
\end{center}

Merge mode shares both motion vector and reference index. It replaces raw MV transmission with small candidate index.

\subsection{Beyond Translational Motion}

Modern codecs model more than 2D translation:

\begin{itemize}
  \item \textbf{affine motion}: rotation, zoom, shear with 4 or 6 parameters;
  \item \textbf{compound prediction}: AV1 can combine Inter and Intra prediction in same block.
\end{itemize}

\begin{center}\includegraphics[width=0.867\linewidth]{ch09/av1-affine-prediction.png}\end{center}
Affine prediction captures motion that a single translational vector cannot model well.

\subsection{Transforms}

Residual transform evolution:

\begin{center}\fitwidth{%
\begin{tabular}{|l|l|}
\hline
\textbf{Standard} & \textbf{Transform design} \\ \hline
H.264/AVC & $4 \times 4$ integer DCT approximation \\ \hline
H.265/HEVC & transforms up to $32 \times 32$, DST for $4 \times 4$ Intra \\ \hline
H.266/VVC & Multiple Transform Selection, non-square transforms \\ \hline
\end{tabular}}
\end{center}

Integer transforms keep encoder/decoder synchronized and avoid mismatch from floating-point precision.

\subsection{Quantization and CABAC}

Quantization is main lossy step. In H.264/HEVC:

\begin{itemize}
  \item QP increase by 6 doubles quantization step size;
  \item QP increase by 1 changes rate by about $12.5\%$ empirically.
\end{itemize}

\begin{calloutbox}{red!70!black}{QP}
Higher QP means coarser quantization, lower bitrate, and higher distortion. QP is primary rate-control knob.
\end{calloutbox}

Modern standards rely on \textbf{CABAC}:

\begin{itemize}
  \item binarizes syntax elements;
  \item adapts probabilities from spatial context;
  \item gives about 5\%-15\% rate reduction over older VLC/CAVLC.
\end{itemize}

\subsection{In-Loop Filters}

Independent block transform and quantization create \textbf{blocking artifacts} at low bitrate.

\begin{center}\includegraphics[width=0.933\linewidth]{ch09/blocking-artifacts.png}\end{center}
If artifacts enter DPB, future motion compensation uses damaged references. Therefore filters must be inside reconstruction loop and normative.

Filter evolution:

\begin{center}\fitwidth{%
\begin{tabular}{|l|l|}
\hline
\textbf{Codec} & \textbf{Filters} \\ \hline
H.264/AVC & deblocking on $4 \times 4$ edges \\ \hline
H.265/HEVC & lighter deblocking on $8 \times 8$ grid plus SAO \\ \hline
H.266/VVC & deblocking, SAO, ALF, LMCS \\ \hline
AV1 & CDEF, loop restoration, film grain synthesis \\ \hline
\end{tabular}}
\end{center}

\subsection{Slices}

Slices divide pictures to contain errors and allow parallel work. At new slice:

\begin{itemize}
  \item Intra prediction does not cross boundary;
  \item CABAC context is reset;
  \item decoder can resynchronize more easily.
\end{itemize}

\begin{center}\includegraphics[width=0.833\linewidth]{ch09/slices-error-containment.png}\end{center}
Slices trade compression efficiency for robustness and parallelism.

\subsection{VCL, NAL, NALU, OBU}

Starting from H.264, standards separate compression from transport:

\begin{itemize}
  \item \textbf{VCL} (\textit{Video Coding Layer}): prediction, transforms, RDO, slices;
  \item \textbf{NAL} (\textit{Network Abstraction Layer}): encapsulation and metadata for transport/storage.
\end{itemize}

\begin{center}\includegraphics[width=0.7\linewidth]{ch09/vcl-nal-layering.png}\end{center}
VCL produces compressed slices; NAL packages them for RTP/IP, MP4, MPEG-2 TS, or raw streams.

NALU types:

\begin{center}\fitwidth{%
\begin{tabular}{|l|l|}
\hline
\textbf{Unit} & \textbf{Meaning} \\ \hline
VCL NALU & compressed slice data: modes, MVs, coefficients \\ \hline
SPS & Sequence Parameter Set: sequence-level parameters \\ \hline
PPS & Picture Parameter Set: per-picture options \\ \hline
SEI & extra metadata: HDR, subtitles, user data \\ \hline
AUD & optional Access Unit Delimiter \\ \hline
\end{tabular}}
\end{center}

\textbf{Access Unit (AU)}: all NALUs needed to decode one frame.

\begin{center}\includegraphics[width=0.933\linewidth]{ch09/nalu-access-unit.png}\end{center}
AU is frame-level decoding group; NALUs are syntax/transport units inside it.

AV1 uses \textbf{OBUs} (\textit{Open Bitstream Units}) instead of MPEG-style NALUs.

\subsection{Annex-B vs Packetized Format}

\textbf{Annex-B byte-stream} inserts start code before each NALU:

\begin{verbatim}
0x00000001
\end{verbatim}

It works for continuous streams but requires byte scanning and emulation prevention.

\begin{center}\includegraphics[width=0.933\linewidth]{ch09/annex-b-byte-stream.png}\end{center}
Annex-B is good for raw \texttt{.264} and MPEG-TS style continuous delivery.

\textbf{Packetized length-prefixed} format writes a length field before each NALU. Parser can jump in $O(1)$, but if container boundaries are damaged, resynchronization is harder.

\begin{center}\includegraphics[width=0.933\linewidth]{ch09/packetized-length-prefixed.png}\end{center}
Packetized format is software-friendly and common in MP4-like containers.

HEVC/VVC keep both formats and expand NALU header to 2 bytes for Temporal ID and Layer ID. AV1 drops Annex-B start codes and uses length-prefixed OBUs.

\subsection{Network-Friendly Operations}

NAL structure lets intermediate nodes inspect stream without decoding:

\begin{itemize}
  \item MANE can drop non-reference B frames under congestion;
  \item sensitive SPS/PPS can receive stronger protection, such as CRC;
  \item RTP can fragment large NALUs into Fragmentation Units and reassemble them.
\end{itemize}

\begin{center}\includegraphics[width=0.933\linewidth]{ch09/nal-network-inspection.png}\end{center}
NAL headers allow network-aware pruning, protection, and packet fragmentation.

\subsection{Random Access}

Random access points define where decoding may restart cleanly.

\begin{center}\includegraphics[width=\linewidth]{ch09/random-access-idr-cra-bla.png}\end{center}
\begin{center}\fitwidth{%
\begin{tabular}{|l|l|l|}
\hline
\textbf{RAP type} & \textbf{Meaning} & \textbf{Use} \\ \hline
\textbf{IDR} & clears DPB; future frames cannot reference before IDR & clean restart, scene cuts \\ \hline
\textbf{CRA} & Intra picture; leading RASL frames may reference past in normal playback & efficient seeking in HEVC+ \\ \hline
\textbf{BLA} & like CRA but signals broken past references & splicing, ad insertion, transcoding \\ \hline
\end{tabular}}
\end{center}

CRA improves compression compared with frequent IDR because normal playback can keep useful past references.

\subsection{Tiles and WPP}

\textbf{Tiles} group CTUs into independent rectangular regions:

\begin{itemize}
  \item self-contained prediction constraints;
  \item parallel encoding/decoding;
  \item useful for viewport-adaptive 360-degree video.
\end{itemize}

\begin{center}\includegraphics[width=0.867\linewidth]{ch09/tiles-parallelism.png}\end{center}
Tiles enable independent rectangular processing regions.

\begin{center}\includegraphics[width=0.867\linewidth]{ch09/tiles-360-video.png}\end{center}
360-degree streaming sends high quality only for viewport tiles and low quality elsewhere.

\begin{calloutbox}{red!70!black}{Tile bandwidth saving}
Tile-based adaptive 360-degree streaming can reduce bandwidth by about 50\%-60\% without degrading perceived viewport quality.
\end{calloutbox}

\textbf{Wavefront Parallel Processing (WPP)} starts CTU row encoding after two CTUs of previous row are ready.

\begin{center}\includegraphics[width=0.867\linewidth]{ch09/wavefront-parallel-processing.png}\end{center}
WPP preserves spatial prediction contexts while enabling row-level multicore parallelism.

\section{Visual Schemes}

\subsection{Modern Hybrid Encoder Loop}

\begin{center}\footnotesize\itshape Diagram (Mermaid source):\end{center}
\begin{verbatim}
flowchart LR
    A[Input frame] --> B[Mode decision]
    B --> C[Intra prediction]
    B --> D[Inter prediction]
    D --> E[Motion compensation]
    C --> F[Prediction]
    E --> F
    A --> G[Residual]
    F --> G
    G --> H[Transform]
    H --> I[Quantization]
    I --> J[Entropy coding]
    J --> K[Bitstream]
    I --> L[Inverse quantization]
    L --> M[Inverse transform]
    M --> N[Reconstruction]
    F --> N
    N --> O[In-loop filters]
    O --> P[Frame buffer]
    P --> D
\end{verbatim}

Internal decoder loop makes encoder use same references as decoder.

\subsection{Standard Abstraction}

\begin{center}\footnotesize\itshape Diagram (Mermaid source):\end{center}
\begin{verbatim}
flowchart LR
    A[Encoder optimization] --> B[Compliant bitstream]
    B --> C[Standard decoder]
    C --> D[Reconstructed video]
    E[Standard] --> B
    E --> C
\end{verbatim}

Standards constrain bitstream and decoder, not full encoder search algorithm.

\subsection{Transport Mapping}

\begin{center}\footnotesize\itshape Diagram (Mermaid source):\end{center}
\begin{verbatim}
flowchart TD
    A[VCL: slices] --> B[NAL units]
    B --> C{Delivery format}
    C --> D[Annex-B start codes]
    C --> E[Length-prefixed packets]
    D --> F[Broadcast / raw stream]
    E --> G[MP4 / packetized storage]
    B --> H[RTP fragmentation]
\end{verbatim}

NAL separates compression syntax from delivery syntax.

\section{Examples}

\subsection{UVG HoneyBee Comparison}

Source sequence: 1080p, 120 fps, target 2.0 Mbps.

\begin{center}\fitwidth{%
\begin{tabular}{|l|r|r|r|}
\hline
\textbf{Codec} & \textbf{Actual rate} & \textbf{PSNR} & \textbf{Encoding time} \\ \hline
MPEG-2 & 7.06 Mbps & 36.58 dB & 2.18 s \\ \hline
H.264 & 2.00 Mbps & 40.51 dB & 17.06 s \\ \hline
HEVC & 1.98 Mbps & 41.69 dB & 109.45 s \\ \hline
AV1 & 1.80 Mbps & 43.47 dB & 158.66 s \\ \hline
VVC & 1.98 Mbps & 43.81 dB & 105.08 s \\ \hline
\end{tabular}}
\end{center}

\begin{calloutbox}{green!55!black}{Takeaway}
MPEG-2 cannot meet target bitrate on this sequence. H.264 reaches target. HEVC, AV1, and VVC give higher PSNR but require far larger encoding time.
\end{calloutbox}

\subsection{UVG Bosphorus Comparison}

Source sequence: 1080p, 120 fps, VBV target 2.0 Mbps.

\begin{center}\fitwidth{%
\begin{tabular}{|l|r|r|r|}
\hline
\textbf{Codec} & \textbf{Actual rate} & \textbf{PSNR} & \textbf{Encoding time} \\ \hline
MPEG-2 & 8.47 Mbps & 36.14 dB & 2.45 s \\ \hline
H.264 & 2.03 Mbps & 35.58 dB & 19.20 s \\ \hline
HEVC & 1.97 Mbps & 38.82 dB & 143.83 s \\ \hline
AV1 & 1.98 Mbps & 40.69 dB & 147.13 s \\ \hline
VVC & 1.96 Mbps & 41.00 dB & 139.08 s \\ \hline
\end{tabular}}
\end{center}

\begin{calloutbox}{blue!60!black}{Missing visual}
Source references HoneyBee and Bosphorus visual codec-comparison figures, but no corresponding image asset was present in available course image folders. Metric tables preserve exam-relevant comparison.
\end{calloutbox}

\subsection{Partition Evolution Summary}

\begin{center}\fitwidth{%
\begin{tabular}{|l|l|l|l|}
\hline
\textbf{Codec} & \textbf{Largest unit} & \textbf{Partition style} & \textbf{Main tradeoff} \\ \hline
H.264/AVC & $16 \times 16$ macroblock & fixed VBS shapes & simple, limited for HD/UHD \\ \hline
HEVC & $64 \times 64$ CTU & quad-tree & good UHD adaptation \\ \hline
VVC & $128 \times 128$ CTU & QT + binary/ternary MTT & best flexibility, high search cost \\ \hline
AV1 & $128 \times 128$ superblock & 10-way split & flexible royalty-free web codec \\ \hline
\end{tabular}}
\end{center}

\subsection{Final Exam Table}

\begin{center}\fitwidth{%
\begin{tabular}{|l|l|}
\hline
\textbf{Concept} & \textbf{Must remember} \\ \hline
hybrid paradigm & same core architecture from MPEG-1 to VVC \\ \hline
standard scope & bitstream + decoder standardized, encoder free \\ \hline
RDO & every mode/partition decided by $J=D+\lambda R$ \\ \hline
complexity growth & more freedom improves compression but expands search tree \\ \hline
Intra evolution & directions grow from H.264 to HEVC to VVC \\ \hline
MPM & predicts Intra mode to reduce signaling \\ \hline
Inter evolution & multiple references, merge mode, fractional ME, affine motion \\ \hline
transform evolution & larger/flexible transforms and MTS for residuals \\ \hline
QP & main lossy rate-control knob; +6 QP doubles step size \\ \hline
CABAC & adaptive arithmetic coding for syntax elements \\ \hline
in-loop filters & protect future references from blocking/ringing \\ \hline
VCL/NAL & compression separated from transport \\ \hline
Annex-B & start-code byte stream for continuous delivery \\ \hline
packetized format & length-prefixed NALUs for containers \\ \hline
AV1 OBU & AV1 replaces NALUs with length-prefixed OBUs \\ \hline
IDR/CRA/BLA & random access and stream splicing tools \\ \hline
tiles & independent regions for parallelism and 360 streaming \\ \hline
WPP & row-level CTU parallelism with two-CTU offset \\ \hline
\end{tabular}}
\end{center}

\begin{calloutbox}{red!70!black}{Main takeaway}
Modern standards improve compression by adding controlled freedom. Encoder must spend computation to search partitions, modes, vectors, transforms, filters, and syntax choices, while decoder remains standardized and deterministic.
\end{calloutbox}
